% Fichier généré automatiquement par tools/generate_corrections.py.
% Source : DYN/DYN-05-Methode/07_RR3D/07_RR3D.tex
% Statut : complété (3/3 question(s)).
\begin{corrigebox}[Corrigé complété]
\CorrigeQuestion{1}
\begin{marginfigure}
\includegraphics[width=\linewidth]{07_RR3D_01_cor}
\end{marginfigure}
\CorrigeQuestion{2}
On isole 2 et on réalise un théorème du moment dynamique en $B$ (ou $A$) en projection sur $\vi{1}$.

On isole 1+2 et on réalise un théorème du moment dynamique en $A$ en projection sur $\vk{0}$.
\CorrigeQuestion{3}
On choisit les inconnues et les conventions positives de la figure,
        on écrit les relations de conservation/fermeture adaptées, puis on résout le
        système obtenu. Le résultat symbolique doit être homogène et vérifié dans les cas
        limites (entrée nulle, rapport unitaire ou position de référence).
\end{corrigebox}
